\section{Hoehn's theorem for the pentagram}

\subsection{Configuration and directed form}

Let $V_1,\ldots,V_5$ be five affine points in general position, with indices
read modulo $5$.  Define the pentagram intersections
\begin{equation}\label{eq:hoehn-W}
 W_i=(V_{i-1}V_{i+1})\cap(V_iV_{i+2}).
\end{equation}
Thus both $W_i$ and $W_{i+1}$ lie on the same diagonal $V_iV_{i+2}$.  Assume
all these intersections are finite and distinct from the relevant endpoints.
Hoehn's two identities, in a uniform directed notation, are
\begin{align}
 \prod_{i=1}^{5}
 \frac{V_iW_i}{W_{i+1}V_{i+2}}&=1,
 \label{eq:hoehn-one}\\
 \prod_{i=1}^{5}
 \frac{V_iW_{i+1}}{W_iV_{i+2}}&=1.
 \label{eq:hoehn-two}
\end{align}
For a convex pentagon, all ten factors are positive, so these are exactly the
ordinary-length formulas in the usual statement.

\subsection{Parameters on the five diagonals}

Write affine homogeneous representatives as
$\widetilde V_i=(x_i,y_i,1)$ and set
\[
 \Delta_{ijk}:=\det(\widetilde V_i,\widetilde V_j,\widetilde V_k).
\]
On the oriented diagonal $V_iV_{i+2}$ define parameters $a_i,b_i$ by
\[
 W_i=(1-a_i)V_i+a_iV_{i+2},
 \qquad
 W_{i+1}=(1-b_i)V_i+b_iV_{i+2}.
\]
Then
\[
 \frac{V_iW_i}{W_{i+1}V_{i+2}}=\frac{a_i}{1-b_i},
 \qquad
 \frac{V_iW_{i+1}}{W_iV_{i+2}}=\frac{b_i}{1-a_i}.
\]

The condition $W_i\in V_{i-1}V_{i+1}$ is
\[
 (1-a_i)\Delta_{i-1,i+1,i}
 +a_i\Delta_{i-1,i+1,i+2}=0.
\]
Introduce
\begin{align*}
 A_i&=\Delta_{i-1,i+1,i},
 &B_i&=-\Delta_{i-1,i+1,i+2},
 &S_i&=A_i+B_i,\\
 C_i&=\Delta_{i+1,i+3,i},
 &D_i&=-\Delta_{i+1,i+3,i+2},
 &T_i&=C_i+D_i.
\end{align*}
The nondegeneracy hypotheses say, in particular, that the denominators below
are nonzero.  Solving the two linear equations gives
\begin{equation}\label{eq:hoehn-parameters}
 a_i=\frac{A_i}{S_i},
 \quad 1-a_i=\frac{B_i}{S_i},
 \qquad
 b_i=\frac{C_i}{T_i},
 \quad 1-b_i=\frac{D_i}{T_i}.
\end{equation}
These formulas can be read entirely cross-multiplied.

\subsection{The cyclic cancellations}

There are three elementary identities:
\begin{equation}\label{eq:hoehn-cyclic}
 D_i=-A_{i+2},
 \qquad C_i=-B_{i-1},
 \qquad T_i=-S_{i+1}.
\end{equation}
The first two are only relabellings, using the alternating and cyclic
symmetries of a determinant.  The third is the four-point affine identity
\begin{multline*}
 \Delta_{i+1,i+3,i}-\Delta_{i+1,i+3,i+2}\\
 +\Delta_{i,i+2,i+1}-\Delta_{i,i+2,i+3}=0.
\end{multline*}
It follows by expanding each determinant, or more conceptually by using the
affine relation among four vectors with final coordinate $1$.  It is a
six-term polynomial identity and requires no area ordering.

Now use \eqref{eq:hoehn-parameters} and \eqref{eq:hoehn-cyclic}:
\begin{align*}
 \frac{a_i}{1-b_i}
 &=\frac{A_iT_i}{S_iD_i}
 =\frac{A_iS_{i+1}}{S_iA_{i+2}},\\
 \frac{b_i}{1-a_i}
 &=\frac{C_iS_i}{T_iB_i}
 =\frac{B_{i-1}S_i}{S_{i+1}B_i}.
\end{align*}
On multiplying for $i=1,\ldots,5$, every $S_i$, every $A_i$, and every $B_i$
appears once in the numerator and once in the denominator.  Both products are
therefore $1$, proving \eqref{eq:hoehn-one} and \eqref{eq:hoehn-two}.

\subsection{A quotient-free certificate}

The two conclusions can be certified without field division as
\begin{align*}
 \prod_i A_iT_i&=\prod_i S_iD_i,\\
 \prod_i C_iS_i&=\prod_i T_iB_i.
\end{align*}
Substitution of \eqref{eq:hoehn-cyclic} turns each side into the same product,
up to ten minus signs.  The ordinary ratio identities follow after cancelling
the explicitly nonzero factors.  Consequently Hoehn's theorem is a fixed-size
ring certificate plus nonzero bookkeeping, and is particularly light for
$\ipra$.
